\thispagestyle{empty}
\begin{center}
\vspace*{26mm}
{\fontsize{26}{30}\selectfont\addfontfeatures{LetterSpace=10}SIDEREVS NVNCIVS}\\[4mm]
{\fontsize{20}{24}\selectfont 星\,际\,信\,使}\\[10mm]
{\large Galilaeus Galilaeus · 伽利略}\\[2mm]
{\small Venetiis, MDCX · 威尼斯，一六一〇年}\\[20mm]
{\color{edgray}\small 拉丁文—中文对照本}\\[6mm]
{\small metachow\ 编译}\\[1mm]
{\color{edgray}\footnotesize Claude（Anthropic）协作翻译、校对与排版}
\end{center}

\vfill
\noindent\begin{minipage}{\textwidth}
\small\color{edgray}\setlength{\parindent}{2em}\linespread{1.25}\selectfont\sloppy
\textbf{编辑说明}\quad
拉丁文以 A. Favaro 编《伽利略全集》国家版第三卷（1892）为底本，并与 1610 年威尼斯初版逐页核对；所据初版为 Smithsonian Libraries 藏本（Internet Archive 数字化）。转写采用现代拼写：æ、œ 拆写为 ae、oe；长 s（ſ）一律作 s；u/v 按音值区分，i 保持不变（\textit{Iuppiter, Ianuarii}）；缩写一律展开；拼法统一为古典形式（\textit{quattuor, numquam, intellego, exspecto}）；普通名词改为小写，天体名称（\textit{Sol, Luna, Tellus, Iuppiter}）保留大写。

本版每页对应初版一页，页眉照录初版页眉，外侧方括号内为初版叶码（如 [7r] 指第 7 叶正面）。拉丁文严格按初版断页，跨页断开的单词归入其开头所在之页。中文尽量在相应分句处断开；中文语序无法拆分时，整句归入其开头所在之页。

字体照初版的分工：初版的罗马体排作正体，中文用宋体；初版的意大利斜体排作斜体，中文用楷体；初版以大写字母强调的词排作字距略开的大写，中文用粗体。

中文译文直接译自拉丁文。插图中，望远镜光路图、几何示意图、星图和木星卫星图按初版逐点测量，重绘为矢量图；月面铜版画、木刻首字母、花饰与印刷商标取自初版扫描，经去黄和折痕修补。

译文由 AI 协作完成，尚未经拉丁文专家审校，错误之处敬请指正。本版译文、注释与排版以 CC BY 4.0 许可发布；拉丁原文与初版扫描属公有领域。源文件：github.com/metachow/sidereus-nuncius-zh
\end{minipage}
\cleardoublepage
